\section{Introduction}
\label{sec:introduction}

Public-safety information systems must remain operational when commercial connectivity fails.
Rescue-service vehicles operate up to several hours away from regional dispatch centres. In that terrain, wireless connectivity alternates between full wide-area coverage, mesh-only operation, and total isolation~\cite{zimbelman2022mesh, wang2023overview}.
Disasters routinely degrade or destroy commercial telecom infrastructure precisely when
coordination demand is highest. In command-hierarchical environments the consequences of
information loss are wrong assignments, missed hazards, and conflicting orders.

The operational constraints are grounded in public Swedish rescue-service
guidance~\cite{rsg-rad201, rsg-kommunalplan}, which treats incident-planning information
as an ongoing, secrecy-restricted process rather than a static artifact; independent
studies of Swedish rescue-service information systems confirm these requirements and the
gaps in information flow~\cite{pilemalm2014enabling, westman2021mobilisation}.

Command at an incident belongs to a \emph{person}, the incident commander, appointed by
the organization. Every resource assignment, status transition, and command hazard
annotation---a \emph{decision}---carries that person's authority. The commander works
from a vehicle (Figure~\ref{fig:setting}), and that vehicle is often cut off from the
regional core. While it is cut off, the organization may decide that command passes to a
deputy or to a superior officer, usually agreed over voice radio. The information system
then holds decisions from two people, and it must say which of them are authoritative.

Ordering is not the hard part. One person cannot be in two vehicles at once, so a
per-person sequence number and GNSS time order that person's decisions, even across a
change of vehicle; epochs and leases are not needed for that case. The hard part is the
cut-over between two people's authority. Clocks give an order, not legitimacy: sorting by
time yields a tidy record of conflicting orders, such as ``engine~3 north'' from the cut-off
commander and ``engine~3 south'' from the deputy, without saying which one was valid.

\begin{figure}[t]
\centering
% The setting AAS operates in: incident site, regional core, intermittent links,
% who writes into the journal, and the command role that can move.
% Hand-drawn TikZ so it shares fonts and colours with Fig. 3: accepting path and
% messages blue (accept), refusal orange (fence). Link state is encoded by line
% style as well as by label: solid = local link, dashed = intermittent radio,
% dotted with a cross = link down.
\begin{tikzpicture}[
  x=1cm, y=1cm,
  font=\scriptsize,
  box/.style={draw=black!55, rounded corners=1.2pt, fill=white, line width=0.5pt,
              align=center, inner sep=2pt},
  cmd/.style={draw=accept, rounded corners=1.2pt, fill=accept!10, line width=1.1pt},
  dim/.style={draw=black!30, rounded corners=1.2pt, fill=white, line width=0.5pt,
              align=center, inner sep=2pt, text=black!50},
  cell/.style={draw=black!55, fill=white, minimum width=0.3cm, minimum height=0.24cm,
               inner sep=0pt, font=\tiny},
  msg/.style={-{Latex[length=1.5mm,width=1.2mm]}, draw=accept, line width=0.6pt},
  local/.style={draw=black!45, line width=0.5pt},
  radio/.style={draw=black!45, line width=0.6pt, dash pattern=on 2.2pt off 1.6pt},
  down/.style={draw=black!25, line width=0.6pt, densely dotted},
  move/.style={-{Latex[length=1.8mm,width=1.5mm]}, draw=black!70, line width=0.8pt},
]
% Incident site: everything a rescue operation brings to the scene.
\draw[rounded corners=3pt, fill=black!4, draw=black!20] (0,0) rectangle (8.6,3.0);
\node[anchor=south west, font=\scriptsize\itshape, text=black!55] at (1.72,0.05) {incident site};

% Regional core, hours away, holding a copy of the journal: a prefix of it.
\draw[rounded corners=1.2pt, fill=white, draw=black!55, line width=0.5pt] (2.4,3.4) rectangle (6.2,4.22);
\node at (4.3,4.03) {\textbf{regional core}, hours away};
\node[anchor=east] at (4.02,3.66) {journal copy};
\foreach \i/\x in {1/4.25, 2/4.55, 3/4.85} {
  \node[cell] at (\x,3.66) {\i};
}

% Command edge: the only sequencer of the incident journal.
\draw[cmd] (3.0,1.35) rectangle (5.6,2.25);
\node at (4.3,2.05) {\textbf{command edge}};
\node[anchor=east] at (4.02,1.62) {journal};
\foreach \i/\x in {1/4.25, 2/4.55, 3/4.85, 4/5.15} {
  \node[cell] at (\x,1.62) {\i};
}

% Backhaul between command edge and core: journal up, epoch and lease down.
\draw[radio] (4.3,2.25) -- (4.3,3.4);
\draw[msg] (4.15,2.3) -- (4.15,3.35);
\draw[msg] (4.45,3.35) -- (4.45,2.3);
\node[anchor=west, text=black!60] at (4.55,2.62) {backhaul};
\node[anchor=east] at (4.08,3.2) {journal};
\node[anchor=west] at (4.52,3.2) {epoch, lease};

% The commander issues decisions at the command vehicle.
\node[box] (commander) at (4.3,0.42) {commander};
\draw[msg] (commander.north) -- (4.3,1.35);
\node[anchor=west] at (4.36,0.98) {decisions};

% Responder edge that is reachable over the mesh.
\node[box, minimum width=1.55cm, minimum height=0.9cm] (r1) at (0.88,1.8) {responder edge\\outbox};
\draw[radio] (r1.east) -- (3.0,1.8);
\draw[msg] (2.55,1.8) -- (2.97,1.8);
\node[anchor=south, text=black!60] at (2.33,1.84) {mesh};
\node[anchor=north] at (2.33,1.76) {submissions};
\node[box] (t1) at (0.88,0.42) {tablet};
\draw[local] (t1.north) -- (r1.south);
\node[anchor=west, text=black!60] at (0.94,0.98) {local};

% Responder edge that is cut off: it keeps working locally, its outbox grows.
\node[dim, minimum width=1.65cm, minimum height=0.9cm] (r2) at (7.65,1.8) {responder edge\\outbox grows};
\draw[down] (5.6,1.8) -- (r2.west);
\node[text=black!60, font=\small] at (6.21,1.8) {$\times$};
\node[anchor=north, text=black!55] at (6.21,1.68) {cut off};
\node[dim] (t2) at (7.65,0.42) {tablet};
\draw[local] (t2.north) -- (r2.south);
\node[anchor=west, text=black!45] at (7.71,0.98) {local};

% The command role can move to another vehicle, by authorized promotion only.
\draw[move] (3.2,2.27) to[out=140, in=40] (1.03,2.27);
\node[anchor=south] at (2.15,2.62) {command role moves};
\end{tikzpicture}

\caption{The setting. The commander's decisions enter the incident journal at the command
vehicle; the regional core, hours away, holds a prefix of it. Dashed links are
intermittent; the responder on the right is cut off but keeps working locally. The
vehicle exercises the commander's authority but holds none of its own. Command passes to
another person only by an organizational decision, and the journal then moves with
them.}
\label{fig:setting}
\end{figure}

Existing designs let something other than the organization choose who may decide.
Consensus elects a leader from whichever quorum is reachable~\cite{ongaro2014raft}, and a
multi-writer conflict-free replicated data type (CRDT) lets every replica write and lets
its merge rule choose between conflicting decisions~\cite{hanssen2025emergency}. Our own
earlier design, \methodname{} (\systemname{}), had the same flaw in a milder form: the
vehicle that obtained an epoch from the core held authority, so reachability still chose
the commander. This paper removes that. The network never decides who is in command; the
system represents the organization's decision, enforces it where it can, and audits it
where it cannot.

Designated authority has a price under partition. While the commander is cut off, a
design can give at most two of three things: the commander keeps recording, command moves
without the commander, and only one disconnected edge accepts decisions
(\S~\ref{sec:tradeoff}). No mechanism escapes this, so succession is a choice of policy,
and each policy has consequences that the organization should be able to see. This leads
to our research question:

\begin{quote}
\textit{How can an intermittently connected incident information system represent,
enforce where possible, and audit succession of command authority from a primary incident
commander to a designated deputy, and what does each succession policy cost when the
commander is cut off?}
\end{quote}

The machinery we use---single-writer journals, outboxes, idempotency keys, epochs, fencing
tokens and leases---is established~\cite{kleppmann2017ddia,gray1989leases}. The
contribution is the authority model and what it exposes:

\begin{enumerate}
  \item (C1) an authority and succession model for incident command in which authority
  belongs to people through signed grants, nodes hold none, a decision binds when it is
  admitted under a valid grant, and decisions made under a superseded grant are kept as
  labelled residuals (\S~\ref{sec:system-model});

  \item (C2) the partition trade-off this model exposes, and a space of six succession
  policies placed in it, each with what it gives up and what it costs
  (\S~\ref{sec:tradeoff}--\S~\ref{sec:policies});

  \item (C3) an evaluation of \systemname{}, the prototype that realizes one of these
  policies, on two independent vehicle edges under a real partition, with TLA+~\cite{tlaplus}
  model checking and baselines against Raft and a last-writer-wins CRDT, together with an
  honest account of where the prototype falls short of the model
  (\S~\ref{sec:architecture}--\S~\ref{sec:evaluation})\footnote{The prototype repository
  is provided as an artifact.}. \pending{Planned: the same scenarios run under each
  policy, on a multi-host testbed.}
\end{enumerate}
